\section{Strategische Ma{\ss}nahmen und ihre Bilanz}
\label{sec:strategie}

\subsection{Mission 30}

Auf dem Kapitalmarkttag im Oktober 2024 legte GEA die Strategie \enquote{Mission 30} mit
drei mittelfristigen Finanzzielen bis 2030 vor: organisches Umsatzwachstum von im Mittel
mehr als \Pz{\MissionWachstum} p.\,a., eine EBITDA-Marge von \Pz{\MissionMargeVon} bis
\Pz{\MissionMargeBis} und ein ROCE von \"uber \Pz{\MissionRoce}\Q{GeaGB}{23}. Gemessen am
Stand 2025 (\Pz{\EbitdaMarge} Marge, \Pz{\Roce} ROCE) liegen beide Gr\"o{\ss}en unter
ihrem Ziel: die Marge knapp unter der Untergrenze der Zielspanne, der ROCE noch deutlich
darunter.

F\"ur 2026 erwartet GEA ein organisches Umsatzwachstum von \Pz{\PrognoseWachstumVon} bis
\Pz{\PrognoseWachstumBis} und eine EBITDA-Marge von \Pz{\PrognoseMargeVon} bis
\Pz{\PrognoseMargeBis}\Q{GeaGB}{2}; nach dem ersten Halbjahr wurde die Erwartung auf
\Pz{\HjWachstumVon} bis \Pz{\HjWachstumBis} beim Wachstum und \Pz{\HjMargeVon} bis
\Pz{\HjMargeBis} bei der Marge angehoben\Q{GeaHJ}{7}.

\subsection{Prognosetreue f\"unf Jahre}

Der Wert einer Strategie h\"angt daran, ob ihre Urheber treffen, was sie ank\"undigen.
Tabelle~\ref{tab:prognose} h\"alt deshalb die Prognose jedes Jahres \emph{in ihrer ersten
Fassung} gegen das Ergebnis der eigenen Reihe. Die sp\"ateren Anpassungen bleiben
au{\ss}en vor: Wer nur die zuletzt angepasste Prognose pr\"uft, misst nichts, weil ein
Unternehmen seine Prognose so lange anpassen kann, bis sie zum Ergebnis passt.

\prognosetabelle

Von \PrognoseQuantifiziert{} quantifizierten Prognosen wurden
\PrognoseTreffer{} erreicht oder \"ubertroffen, keine verfehlt. Die qualitativen Angaben
(\enquote{significantly rising}) bleiben unbewertet, weil sie keine Grenze nennen, an der
sie scheitern k\"onnten.

\bk{Ein Ergebnis von \PrognoseTreffer{} aus \PrognoseQuantifiziert{} l\"asst zwei Lesarten
zu. Entweder ist die Planung genau, oder die Prognosen sind vorsichtig gesetzt. Daf\"ur
spricht, dass die Marge in beiden Jahren, in denen sie \"uberhaupt als Spanne vorgegeben
wurde, \"uber der Spanne endete, und dass GEA die Prognose unterj\"ahrig regelm\"a{\ss}ig
anhob. Beide Lesarten sind f\"ur den Anleger gute Nachrichten, aber sie bedeuten
Verschiedenes: Vorsicht in der Planung hei{\ss}t, dass die Ziele von Mission 30 eher eine
Untergrenze als ein Wunsch sind.}

\subsection{Was Mission 30 rechnerisch verlangt}

Die beiden Ziele lassen sich in eine einzige Wachstumsrate \"ubersetzen. Erreicht GEA bis
2030 ein organisches Umsatzwachstum von \Pz{\MissionWachstum} p.\,a.\ und die Mitte der
Margenspanne von \Pz{\MissionMargeMitte}, so w\"achst das EBITDA vor Restrukturierung um
\MissionImplizit{} p.\,a.\ gegen\"uber dem Stand von 2025.

\annahme{Die Rechnung unterstellt, dass die Margenspanne linear bis 2030 erreicht wird und
dass Zuk\"aufe unber\"ucksichtigt bleiben; GEA nennt das Wachstumsziel ausdr\"ucklich
organisch. Sie unterstellt ferner, dass sich das EBITDA-Wachstum in gleichem Ma{\ss}e im
freien Zufluss niederschl\"agt. Die Reihe in Tabelle~\ref{tab:reihe} zeigt, dass genau
das in den vergangenen Jahren nicht geschah: Das EBITDA stieg, der freie Zufluss nicht,
weil die Investitionen st\"arker wuchsen. Die \MissionImplizit{} sind damit eine
Obergrenze, kein Erwartungswert.}

\subsection{Kapitalverwendung}

Das zwischen November 2023 und April 2025 durchgef\"uhrte R\"uckkaufprogramm erwarb
\Mio{\RueckkaufStueckAlt}{Aktien} Aktien zu durchschnittlich \je{\RueckkaufPreisAlt}{EUR}, ein
Volumen von \Mio{\RueckkaufVolumenAlt}{EUR}\Q{GeaGB}{14}. Ein neues Programm \"uber
\Mio{\RueckkaufProgramm}{EUR} l\"auft; bis zum Halbjahresstichtag waren
\Mio{\RueckkaufStueck}{Aktien} Aktien erworben, und eine zweite Tranche \"uber
\Mio{\RueckkaufTrancheZwei}{EUR} wurde angek\"undigt\Q{GeaHJ}{44}.

\bk{Die R\"uckk\"aufe des abgeschlossenen Programms erfolgten zu einem Durchschnittskurs
von \je{\RueckkaufPreisAlt}{EUR}, also weit unter dem heutigen Kurs. Das war
Kapitalverwendung zu einem guten Preis. Das laufende Programm kauft zu deutlich
h\"oheren Kursen; dieselbe Ma{\ss}nahme schafft damit heute weniger Wert je eingesetztem
Euro als vor zwei Jahren. Ein R\"uckkauf ist eine Investition in die eigene Aktie und
unterliegt derselben Preisfrage wie jede andere.}
